\subsection{零维空间与卢津空间}

定义 5. 拓扑空间称为零维的，如果它是豪斯多夫的，且每个点有一个既开又闭的邻域基本系。

每个零维空间 X 都是全不连通的；因为一点的连通分支含于所有包含该点且既开又闭的集（第 I 章，§ 11，第 5 号），而这些集的交正是 \(\{x\}\) ，若 \(\mathbf{X}\) 是零维的。

反之，局部紧全不连通空间是零维的（第 II 章，§ 4，第 4 号，命题 6）；但存在不是零维的全不连通可度量化空间{[}习题 3 b){]}。

零维空间的每个子空间是零维的，且零维空间的拓扑和与积是零维的。

定义 6. 拓扑空间 X 是卢津空间，如果它可度量化，且存在一个零维波兰空间 P 和 P 到 X 上的连续双射映射。

显然每个卢津空间都是苏斯林空间。

命题 12. 可度量化空间是卢津空间，当且仅当它是一个波兰空间在连续双射映射下的像。

此条件显然必要；我们证明它也充分。若 f 是卢津空间 X 到可度量化空间 Y 上的连续双射，则由定义 6 ，Y 是卢津空间。故只需证明波兰空间是卢津空间。

首先注意，若 X 是卢津空间，则 X 的每个闭（相应地，开）子空间 A 是卢津空间（参看第 7 号，定理 3）；因为若 f 是零维波兰空间 P 到 X 上的连续双射，则 \(\overline{f}^{1}(A)\) 在 P 中闭（相应地，开），因而是 P 的零维波兰子空间（第 1 号，命题 1 与命题 2）。

卢津空间的每个可数积是卢津空间；这由第 1 号的命题 1 及零维空间的每个积是零维的这一事实得出。豪斯多夫拓扑空间中卢津子空间的每个可数交是卢津子空间；这由前面的注及第 1 号的引理 1 得出。此外：

引理 2. 若可度量化空间 X 使得存在 X 的一个由卢津子空间组成的可数分划 \((\mathbf{A}_{n})\) ，则 X 是卢津空间。

对每个 n ，设 \(P_{n}\) 是零维波兰空间，\(f_{n}\) 是 \(P_{n}\) 到 \(A_{n}\) 上的连续双射。若 P 是诸 \(P_{n}\) 的拓扑和，则 P 是波兰的（第 1 号，命题 1）且零维，而映射 \(f: P \to X\) （它与 \(f_{n}\) 在每个 \(P_{n}\) 上（对每个 n ）一致）是 P 到 X 上的连续双射；由此即得结论。

现在证明 R 的区间 \(I = [0, i]\) 是卢津空间。设 J 是由 I 中所有无理数组成的子空间；则 J 是波兰空间（第 1 号，命题 3 的推论）；且 J 是零维的，因为若 x 是 J 的任意点，则 R 中形如 {]}r, s{[} （其中 r 与 s 是有理数且 r \textless{} x \textless{} s ）的区间在 J 上的迹构成 x 在 J 中的既开又闭邻域的基本系（因为 {]}r, s{[} 与 {[}r, s{]} 在 J 上的迹相同）。故 J 是 I 的卢津子空间。而 J 与 I 的每个由单个有理点组成的子空间构成 I 的可数分划，因此由引理 2 ，I 是卢津空间。

设 P 是任意波兰空间。由定理 1 的推论 1 （第 1 号），P 同胚于方体 IN 的一个子空间，后者是 IN 中开集的可数交。由于 I 是卢津空间，证明开头的诸注表明 P 是卢津空间，从而命题 12 的证明完成。

